% Minimal retained preamble for lane-p2 cards taken from
% sources/sigma-2010-081-generalized-complex-reduction/pinned-original-source.tex
% (Singular Reduction of Generalized Complex Manifolds, DOI 10.3842/SIGMA.2010.081).
%
% Every declaration below is copied verbatim from the author's own file:
%   lines 9-19  : the author's theorem declarations
%   line  21    : \numberwithin{equation}{section}
%   line  23    : \usepackage[all]{xy}
%   lines 26-36 : the author's own macro definitions
% The published paper uses the publisher class \documentclass{sigma} (chosen in
% the \ifpdf block at lines 1-6), which the assembler does not accept -- only
% article and amsart are supported, see NEW_TASK_AUTHORING section 16.1.  The
% wrapper therefore supplies the article class and the base packages
% amsmath/amssymb/amsthm/graphicx.  Nothing was renamed, redefined, added or
% removed; no macro was invented.  xy is retained because the retained body
% contains the author's own commutative diagram.
\documentclass{article}
\usepackage{amsmath,amssymb,amsthm,graphicx}
\usepackage[all]{xy}

\newtheorem{Theorem}{Theorem}[section]
\newtheorem{Lemma}[Theorem]{Lemma}
%\newtheorem{problem}[Theorem]{Problem}
%\newtheorem{conv}[Theorem]{Convention}
\newtheorem{Proposition}[Theorem]{Proposition}
\newtheorem{Corollary}[Theorem]{Corollary}
{\theoremstyle{definition}
\newtheorem{Remark}[Theorem]{Remark}
\newtheorem{Definition}[Theorem]{Definition}
\newtheorem{Example}[Theorem]{Example}
}

\numberwithin{equation}{section}

\newcommand{\deff}[1]{\textbf{\emph{#1}}}
\newcommand{\func}[3]{#1 \colon #2 \to #3}
\newcommand{\hookfunc}[3]{#1 \colon #2 \hookrightarrow #3}
\newcommand{\eval}[2]{\left. #1 \right|_{#2}}
\newcommand{\pont}[1]{\mathbb{T} {#1}}
\newcommand{\cpont}[1]{\mathbb{T}_{\mathbb{C}} {#1}}
%\newcommand{\tang}{\pmb{\top}}
\newcommand{\tang}{T}
\newcommand{\iprod}[2]{\left\langle {#1} , {#2} \right\rangle}
\newcommand{\iiprod}[2]{\left\langle \! \left\langle {#1} , {#2} \right\rangle \! \right\rangle}
\newcommand{\litlet}{\renewcommand{\labelenumi}{(\alph{enumi})}}
\title{A Hamiltonian action restricts to a G-stable twisted generalized complex submanifold}
\author{Timothy E. Goldberg}
\date{doi:10.3842/SIGMA.2010.081; CC BY 4.0}
\begin{document}
\maketitle
\noindent\textit{Editorial note.} The counted unit is the article's own Theorem thm6 at source lines 617-620 together with the proof environment that begins on the next line after the blank lines following the statement (623-666); the pairing is the ordinary adjacent one.\par\medskip
\noindent\textit{Editorial note.} The retained context at source lines 187-193 is the article's own definition of GC structures, twisted GC structures and their manifolds; the retained context at source lines 302-304 is its own definition of a GC submanifold and of the induced structure E\_S; the retained context at source lines 570-588 is its own definition of generalized Hamiltonian and twisted generalized Hamiltonian actions together with the generalized moment map and the moment one-form.\par\medskip
\noindent\textit{Editorial note.} Licence: the licence evidence is not a JSON observation file for this source. The card relies on the article-specific attribution notice retained with the corpus at editable-corpus/licenses/2819-source-license.txt, which names the article, the DOI, the publisher and CC BY 4.0, together with the retained CC BY 4.0 legal code at editable-corpus/licenses/CC-BY-4.0.txt; the paths are given relative to the corpus root. No acceptance record was written and no new observation was made.\par\medskip
\noindent\textit{Editorial note.} The published paper is typeset with the publisher class sigma, which the assembler does not accept (article and amsart only). The retained preamble supplied with this task therefore reproduces, verbatim, the author's theorem declarations at source lines 9-19, his numberwithin line at source line 21, his xy package line at source line 23 and his macro definitions at source lines 26-36, wrapped in the supported article class with amsmath, amssymb, amsthm and graphicx. Nothing was renamed, redefined or added.\par\medskip
\noindent\textit{Editorial note.} No retained line contains a citation or a cross-reference, so the material carries no bibliography and needs no navigation pointer.\par\medskip
\noindent\textit{Editorial note.} The article's main theorem MainThm, its corollary MainThm2, its Proposition prop9 and its other propositions, lemmas and examples are not counted here and are not used as premises.\par\medskip
\noindent\textit{Allowed standard prerequisites (editorial).}\begin{itemize}
\item Twisted GC structure and twisted GC manifold. An almost GC structure is a GC structure when the corresponding complex Dirac structure is Courant involutive; given a closed three-form H, it is an H-twisted GC structure when the structure is H-twisted Courant involutive, and (M,E,H) is then a twisted GC manifold; retained at source lines 187-193.. Fixes the three objects of the hypothesis of the counted theorem, namely the twisted GC manifold (M,E,H) with its associated complex Dirac structure E.
\item Twisted GC submanifold and the induced structure E\_S. For a submanifold S of a twisted GC manifold (M,E,H), if E\_S is a subbundle of the complexified generalized tangent bundle over S meeting its conjugate trivially, then (S,E\_S,j\textasciicircum{}*H) is a (twisted) GC submanifold of (M,E,H), with J\_S the GC structure induced by E\_S; retained at source lines 302-304.. Supplies the object S of the counted theorem and the structure E\_S whose G-stability the proof establishes first.
\item Twisted generalized Hamiltonian action, moment map and moment one-form. For a Lie group G acting canonically on a twisted GC manifold (M,J,H), the action is twisted generalized Hamiltonian when there are a G-equivariant map mu into the dual of the Lie algebra and a G-equivariant one-form alpha with values in that dual such that for every xi one has xi\_M = -J(d mu\textasciicircum{}xi) - alpha\textasciicircum{}xi, equivalently xi\_M + alpha\textasciicircum{}xi - i d mu\textasciicircum{}xi a section of E, and iota\_\{xi\_M\} H = d alpha\textasciicircum{}xi; mu and alpha are called a generalized moment map and a moment one-form; retained at source lines 570-588.. Gives the two defining identities that the proof has to reproduce on S, with mu restricted and alpha pulled back.
\end{itemize}\medskip
\begin{Definition}
Let $M$ be a manifold, and let $L$ be a real (respectively complex) linear subbundle of $\pont{M}$ (respectively $\cpont{M}$).  Then $L$ is \deff{Courant involutive} if the space $\Gamma(L)$ of smooth sections of $L$ is closed under the Courant bracket, i.e. $\big[ \Gamma(L), \Gamma(L) \big] \subset \Gamma(L)$.  If $H \in \Omega_{\text{cl}}^3(M)$, we similarly def\/ine \deff{$H$-twisted Courant involutive}.

Let $E \subset \cpont{M}$ be an almost GC structure on $M$.  This is a \deff{GC structure} if $E$ is Courant involutive, in which case $(M,E)$ is a \deff{GC manifold}.  If $H \in \Omega_{\text{cl}}^3(M)$ and $E$ is $H$-twisted Courant involutive, then $E$ is an \deff{$H$-twisted GC structure}, and $(M,E,H)$ is a \deff{twisted GC manifold}.

Let $(\mathcal{J}_1, \mathcal{J}_2)$ be an almost GK structure on $M$.  This is a \deff{GK structure} if both $\mathcal{J}_1$ and~$\mathcal{J}_2$ are Courant involutive, in which case $(M,\mathcal{J}_1,\mathcal{J}_2)$ is a \deff{GK manifold}.  If $H \in \Omega_{\text{cl}}^3(M)$ and~$\mathcal{J}_1$ and $\mathcal{J}_2$ are $H$-twisted Courant involutive, then this is an \deff{$H$-twisted GK structure}, and $(M,\mathcal{J}_1,\mathcal{J}_2,H)$ is a \deff{twisted GK manifold}.
\end{Definition}

\begin{Definition}
Let $(M,E,H)$ be a twisted GC manifold, and let $\hookfunc{j}{S}{M}$ be a submanifold.  If $E_S \subset \cpont{S}$ is a subbundle and satisf\/ies $E_S \cap \overline{E_S} = 0$, then $(S,E_S,j^* H)$ is a \deff{(twisted) GC submanifold} of $(M,E,H)$, and we denote by $\mathcal{J}_S$ the GC structure on $S$ induced by $E_S$.
\end{Definition}

\begin{Definition}
\label{genhamactiondef}

Let $(M,\mathcal{J})$ be an untwisted GC manifold, let $E$ be the associated complex Dirac structure on $M$, and let $G$ be a Lie group acting canonically on $(M,\mathcal{J})$.  This action is \deff{generalized Hamiltonian} if there exists a $G$-equivariant map $\func{\mu}{M}{\mathfrak{g}^{\ast}}$ such that, for all $\xi \in \mathfrak{g}$,
\[
  \xi_M = - \mathcal{J} ( \mathrm{d} \mu^{\xi} )
\]
or equivalently $\xi_M - i \mathrm{d} \mu^{\xi} \in \Gamma(E)$.
Here $\func{\mu^{\xi}}{M}{\mathbb{R}}$ is the smooth function def\/ined by $\mu^{\xi}(x) := \iprod{\mu(x)}{\xi}$ for all $x \in M$.  The map $\mu$ is called a \deff{generalized moment map} for the $G$-action on~$(M, \mathcal{J})$.



Let $(M,\mathcal{J},H)$ be a twisted GC manifold, and let $G$ be a Lie group acting canonically on $(M,\mathcal{J},H)$.  This action is \deff{twisted generalized Hamiltonian} if there exists a $G$-equivariant map $\func{\mu}{M}{\mathfrak{g}^{\ast}}$ and a $G$-equivariant $\mathfrak{g}^*$-valued one-form $\alpha \in \Omega^1(M, \mathfrak{g}^{\ast})$ on $M$ such that, for all $\xi \in M$,
\begin{enumerate}\itemsep=0pt
  \item $\xi_M = - \mathcal{J} ( \mathrm{d} \mu^{\xi} ) - \alpha^{\xi}$ (or equivalently $\xi_M + \alpha^{\xi} - i \mathrm{d} \mu^{\xi} \in \Gamma(E)$), and
  \item $\iota_{\xi_M} H = \mathrm{d} \alpha^{\xi}$.
\end{enumerate}
Here $\mu^{\xi}$ is as def\/ined above, and $\alpha^{\xi} \in \Omega^1(M)$ is the dif\/ferential one-form on $M$ def\/ined by $(\alpha^{\xi})_x (v) := \iprod{\alpha_x(v)}{\xi}$ for all $x \in M$, $v \in \tang_x M$.  The map $\mu$ and the one-form $\alpha$ are called a~\deff{generalized moment map} and a \deff{moment one-form}, respectively, for the action $G$-action on $(M, \mathcal{J}, H)$.
\end{Definition}

\begin{Theorem} \label{thm6}
Let $(M,E,H)$ be a twisted GC manifold, where $E$ is the associated complex Dirac structure, and let $G$ be a Lie group acting on $(M,E,H)$ in a Hamiltonian fashion with moment map $\func{\mu}{M}{\mathfrak{g}^*}$ and moment one-form $\alpha \in \Omega^1(M,\mathfrak{g}^*)$.
If $\hookfunc{j}{S}{M}$ is a $G$-stable twisted GC submanifold of $(M,E,H)$, then the restriction of the action of $G$ to $(S,E_S,j^* H)$ is Hamiltonian with moment map $\func{\mu|_S}{S}{\mathfrak{g}^*}$ and moment one-form $j^* \alpha \in \Omega^1(S,\mathfrak{g}^*)$.
\end{Theorem}

\begin{proof}
First we will prove that the action of $G$ on $S$ preserves $E_S$.  Let $x \in S$ and $(X, \lambda) \in (\tang_{\mathbb{C}, x} S \oplus \tang_{\mathbb{C}, x}^{\ast} M) \cap E_x$, which means that $(X, \lambda|_S) \in E_{S,x}$.  Then for any $g \in G$ we have
\[
  g \cdot (X + j^* \lambda) = g_* (X) + (g^{-1})^* \lambda|_S.
\]
Because $S$ is $G$-stable, the inclusion $\hookfunc{j}{S}{M}$ is $G$-equivariant, i.e. the $G$-action commutes with $j$.  Hence $\hookfunc{j_*}{\tang S}{\tang M}$ is $G$-equivariant, so $g_* (X) \in \tang_{g \cdot x} S$.  Also
\[
  (g^{-1})^* \left( \lambda|_S \right)
  =
  (g^{-1})^* j^* \lambda
  =
  j^* (g^{-1})^* \lambda
  =
  \left( (g^{-1})^* \lambda \right)|_S.
\]
Since $E$ is $G$-stable, we have $g \cdot (X + \lambda) = g_* (X) + (g^{-1})^* \lambda \in E_{g \cdot x}$.  Therefore $g \cdot (X, j^* \lambda) = \left( g_* (X), j^* (g^{-1})^* \lambda \right) \in E_{S, g \cdot x}$.  Thus $E_S$ is $G$-stable.


Now suppose that $(S,E_S,j^* H)$ is a GC submanifold of $(M,E,H)$, meaning that $E_S$ is a~vector bundle, that $E_S \cap \overline{E_S} = 0$, and that $E_S$ is $j^* H$-twisted Courant involutive.  Since $j$ is $G$-equivariant, for all $\xi \in \mathfrak{g}$ we have $\xi_M|_S = \xi_S$, $(j^* \mu)^{\xi} = j^* (\mu^{\xi})$, and $(j^* \alpha)^{\xi} = j^* (\alpha^{\xi})$.  Furthermore, by the naturality of the exterior derivative we have
\[
  \mathrm{d} (j^* \mu)^{\xi}
  =
  \mathrm{d} j^* (\mu^{\xi})
  =
  j^* \big(\mathrm{d} \mu^{\xi}\big),
\]
so $\mathrm{d} (\mu|_{S})^{\xi} = (\mathrm{d} \mu^{\xi})|_S$.  For each $x \in S \subset M$, since $\eval{\left( \xi_M + \alpha^{\xi} - i \, \mathrm{d} \mu^{\xi} \right)}{x} \in E_x$, this means that
\[
  \eval{\big( \xi_S + \alpha^{\xi}|_S - i \, \mathrm{d} \mu^{\xi}\big|_S \big)}{x} \in E_{S,x}.
\]
Again using the $G$-equivariance of $j$, for all $x \in S$ we have
\[
  \eval{\iota_{\xi_S} j^* (H)}{x}
  =
  \eval{j^* (\iota_{\xi_S} H)}{x}
  =
  \eval{j^* (\iota_{\xi_M} H)}{x}
  =
  \eval{j^* (\mathrm{d} \alpha^{\xi})}{x}
  =
  \eval{\mathrm{d} (j^* \alpha)^{\xi}}{x}.
\]
Thus the action of $G$ on $(S,E_S,j^* H)$ is twisted Hamiltonian with moment map $\mu|_S$ and moment one-form $\alpha|_S$.
\end{proof}

\end{document}
